% Source: physical PDF pp. 279--289 (printed pp. 256--266), from the
% Section 29.2 heading through the final sentence immediately before
% Section 29.3.
% Inventory: Eqs. (29.2.1)--(29.2.51), ten unnumbered display groups,
% three ordinary footnotes, one centered three-asterisk divider, and no
% citations, figures, or tables.
% Corrected the mass argument of \Delta in Eq. (29.2.40).
% Uncertainties: none.

\subsection{Supersymmetry Current Sum Rules}

We now turn to sum rules that yield exact quantitative relations between
the vacuum energy and parameters describing the strength of supersymmetry
breaking.

Let us again start by assuming that the world is placed in a box of volume
\(V\), with periodic boundary conditions to preserve translation
invariance. The vacuum expectation value of the anticommutation relations
\eqref{eq:25.2.36} may then be expressed as a sum over discrete states
\(\ket{X,\mathrm{Box}}\):
\begin{equation}
  \begin{aligned}
  &\sum_X
    \bra{\mathrm{VAC}} Q_\alpha\ket{X,\mathrm{Box}}
    \langle X,\mathrm{Box}\lvert
      \bar Q_{\dot\beta}\rvert\mathrm{VAC}\rangle
  \\
  &\quad
  +\sum_X
    \langle\mathrm{VAC}\lvert\bar Q_{\dot\beta}
      \rvert X,\mathrm{Box}\rangle
    \bra{X,\mathrm{Box}} Q_\alpha\ket{\mathrm{VAC}}
  \\
  &\hspace{30mm}
  =-2\sigma_{\mu\,\alpha\dot\beta}
    \bra{\mathrm{VAC}} P^\mu\ket{\mathrm{VAC}}\,.
  \end{aligned}
  \tag{29.2.1}\label{eq:29.2.1}
\end{equation}
Here the label `Box' indicates that the states are normalized to have
Kronecker deltas rather than delta functions as their scalar products.
Setting \(\beta=\alpha\), summing over \(\alpha\), and using
Eq.~\eqref{eq:25.2.37} gives
\begin{equation}
  \sum_{X,\alpha}
  \left\lvert
    \bra{\mathrm{VAC}} Q_\alpha\ket{X,\mathrm{Box}}
  \right\rvert^2
  =4\bra{\mathrm{VAC}} P^0\ket{\mathrm{VAC}}\,.
  \tag{29.2.2}\label{eq:29.2.2}
\end{equation}

Because \(Q_\alpha\) commutes with the four-momentum, it is only states
with zero three-momentum and the same energy as the vacuum that can
contribute to this sum. To find the volume dependence of the matrix
elements that do contribute in Eq.~\eqref{eq:29.2.2}, we note that a
box-normalized state \(\ket{X,\mathrm{Box}}\) containing \(N_X\)
particles is related according to Eq.~\eqref{eq:3.4.3} to the
corresponding continuum-normalized state \(\ket{X}\) by
\begin{equation}
  \ket{X,\mathrm{Box}}
  =\left((2\pi)^3/V\right)^{N_X/2}\ket{X}\,.
  \tag{29.2.3}\label{eq:29.2.3}
\end{equation}
For states with \(\mathbf p_X=0\), the spatial integral of the
supersymmetry current time-component \(S_\alpha^0\) gives another factor
of \(V\), so we conclude that for box-normalized states with
\(\mathbf p_X=0\)
\begin{equation}
  \bra{\mathrm{VAC}} Q_\alpha\ket{X,\mathrm{Box}}
  =(2\pi)^{3N_X/2}V^{1-N_X/2}
  \bra{\mathrm{VAC}} S_\alpha^0(0)\ket{X}\,.
  \tag{29.2.4}\label{eq:29.2.4}
\end{equation}
Since invariance under \(2\pi\)-rotations does not allow \(X\) to be a
zero-particle state, the dominant terms in Eq.~\eqref{eq:29.2.2} when
\(V\to\infty\) will be those from one-particle states. In this limit,
Eq.~\eqref{eq:29.2.2} becomes
\begin{equation}
  (2\pi)^3\sum_{X,\alpha}^{(0)}
  \left\lvert
    \bra{\mathrm{VAC}} S_\alpha^0(0)\ket{X}
  \right\rvert^2
  =4\rho_{\mathrm{VAC}}\,,
  \tag{29.2.5}\label{eq:29.2.5}
\end{equation}
where \(\rho_{\mathrm{VAC}}\) is the vacuum energy density
\begin{equation}
  \rho_{\mathrm{VAC}}
  \equiv
  \bra{\mathrm{VAC}} P^0\ket{\mathrm{VAC}}/V\,,
  \tag{29.2.6}\label{eq:29.2.6}
\end{equation}
and the superscript \((0)\) indicates that the sum in
Eq.~\eqref{eq:29.2.5} runs only over one-particle states with zero
four-momentum. These of course are the two helicity states of the
goldstino.

We see again from Eq.~\eqref{eq:29.2.5} that if the vacuum energy density
does not vanish then the vacuum is not invariant under supersymmetry
transformations, but is rather transformed into one-goldstino states.
Conversely, Eq.~\eqref{eq:29.2.2} shows that if the vacuum is not invariant
under supersymmetry then according to Eq.~\eqref{eq:29.2.2} its energy in
a finite box cannot vanish, although it is conceivable that supersymmetry
transformations might take the vacuum only into multiparticle states, in
which case the vacuum energy density would vanish in the limit of large
volume.

To evaluate the one-goldstino contribution in Eq.~\eqref{eq:29.2.5}, we
use Lorentz invariance to write the matrix element of the supersymmetry
current between the vacuum and a one-goldstino state
\(\ket{\mathbf p,\lambda}\), with momentum \(\mathbf p\) and
helicity \(\lambda\), in the form\footnote{Lorentz invariance alone would
yield this formula with independent coefficients \(F_L,F'_L\) and
\(F_R,F'_R\) for the undotted and dotted current components,
respectively. It is CPT invariance that imposes the
relations \(F_R=F_L^*\) and \(F'_R=F_L'{}^*\). To see this, we must use the
CPT-transformation properties of the supersymmetry current
\[
  \mathrm{CPT}\,S^\mu_\alpha(x)(\mathrm{CPT})^{-1}
  =-\bar S^\mu_{\dot\alpha}(-x),
  \qquad
  \mathrm{CPT}\,\bar S^\mu_{\dot\alpha}(x)(\mathrm{CPT})^{-1}
  =S^\mu_\alpha(-x)
\]
(see Section 5.8) and of the one-particle states
\[
  \mathrm{CPT}\ket{\mathbf p,\lambda}
  =\chi_\lambda\ket{\mathbf p,-\lambda}
\]
with \(\chi_\lambda\) a phase factor that depends on how we define the
relative phases of the helicity states. We also need the reality
properties of the coefficient functions \(u(\mathbf p,\lambda)\). These
are related to those of the one-particle states by the definition
\[
  \langle\mathrm{VAC}\lvert\psi_{\mathrm{REN},\alpha}(x)
  \rvert\mathbf p,\lambda\rangle
  =(2\pi)^{-3/2}\exp(\ii p\mathbin{\cdot}x)
    u_\alpha(\mathbf p,\lambda)\,,
\]
where \(\psi_{\mathrm{REN},\alpha}(x)\) is a renormalized Weyl field,
with dotted conjugate and CPT-transformation properties
\[
  \mathrm{CPT}\,\psi_{\mathrm{REN},\alpha}(x)(\mathrm{CPT})^{-1}
  =\bar\psi_{\mathrm{REN},\dot\alpha}(-x),
  \qquad
  \mathrm{CPT}\,\bar\psi_{\mathrm{REN},\dot\alpha}(x)
    (\mathrm{CPT})^{-1}
  =-\psi_{\mathrm{REN},\alpha}(-x)\,,
\]
which yields
\(\mathrm{CPT}:u_\alpha(\mathbf p,\lambda)
\mapsto\chi_\lambda^*\bar u_{\dot\alpha}(\mathbf p,-\lambda)\), up to
the displayed phase convention. Putting this together with the
CPT-transformation properties of the supersymmetry current and
one-particle states yields the relations \(F_R=F_L^*\) and
\(F'_R=F_L'{}^*\).}
\begin{equation}
  \begin{aligned}
  \langle\mathrm{VAC}\lvert S^\mu_\alpha(0)
    \rvert\mathbf p,\lambda\rangle
  &=(2\pi)^{-3/2}\left[
    -\ii F\sigma^\mu_{\alpha\dot\beta}
      \bar u^{\dot\beta}(\mathbf p,\lambda)
    +\ii p^\mu F'u_\alpha(\mathbf p,\lambda)\right],
  \\
  \langle\mathrm{VAC}\lvert\bar S^\mu_{\dot\alpha}(0)
    \rvert\mathbf p,\lambda\rangle
  &=(2\pi)^{-3/2}\left[
    -\ii F^*\bar\sigma^\mu_{\dot\alpha\beta}
      u^\beta(\mathbf p,\lambda)
    +\ii p^\mu F'{}^*\bar u_{\dot\alpha}(\mathbf p,\lambda)\right],
  \end{aligned}
  \tag{29.2.7}\label{eq:29.2.7}
\end{equation}
where \(u_\alpha(\mathbf p,\lambda)\) and
\(\bar u_{\dot\alpha}(\mathbf p,\lambda)\) are the two-component
coefficient functions for a massless field, and \(F\) and \(F'\) are unknown
constants. The matrix element \eqref{eq:29.2.7} satisfies a conservation
condition
\(p_\mu\langle\mathrm{VAC}\lvert S^\mu_\alpha(0)
\rvert\mathbf p,\lambda\rangle=0\), and likewise for the dotted current,
because the wave functions satisfy the two Weyl blocks of the
momentum-space equation \eqref{eq:5.5.42} for zero
mass, and \(p^\mu\) is on the light-cone. The sum over helicities gives
\[
  \sum_\lambda
    u_\alpha(\mathbf p,\lambda)\bar u_{\dot\beta}(\mathbf p,\lambda)
  =-\frac{p_\mu\sigma^\mu_{\alpha\dot\beta}}{2p^0},
  \qquad
  \sum_\lambda
    \bar u^{\dot\alpha}(\mathbf p,\lambda)u^\beta(\mathbf p,\lambda)
  =-\frac{p_\mu\bar\sigma^{\mu\dot\alpha\beta}}{2p^0}\,.
\]
(The Weyl wave functions for a massless particle of
momentum \(\mathbf p\) and spin \(\lambda\) are not well defined for
\(\mathbf p\to0\), but there is no problem if we regard the sum over
\(X\) in Eq.~\eqref{eq:29.2.5} as a sum over helicities for a small
momentum \(\mathbf p\) of fixed direction.) After a straightforward
calculation, we then have
\[
  4\rho_{\mathrm{VAC}}
  =-\frac{\lvert F\rvert^2}{2p^0}
    \operatorname{tr}_2\left(
      \sigma^0p_\mu\bar\sigma^\mu
      +\bar\sigma^0p_\mu\sigma^\mu\right)
  =2\lvert F\rvert^2\,,
\]
and therefore
\begin{equation}
  \rho_{\mathrm{VAC}}=\lvert F\rvert^2/2\,.
  \tag{29.2.8}\label{eq:29.2.8}
\end{equation}
An alternative proof of this formula will be given at the end of this
section.

The parameter \(F\) plays much the same role in the interactions of soft
goldstinos as does the parameter \(F_\pi\) (introduced in Section 19.4) in
the interactions of soft pions. The matrix element of the supersymmetry
current between any two states \(X\) and \(Y\) may be split into terms
that have a one-goldstino pole at \(p^\mu=0\) in the momentum transfer
\(p\equiv p_X-p_Y\) and terms that do not have this pole:
\begin{equation}
  \begin{aligned}
  \bra{X} S^\mu_\alpha(0)\ket{Y}
  ={}&\frac{\ii}{p^2}\left[
    F(\sigma^\mu p_\nu\bar\sigma^\nu)_\alpha{}^\beta
      \mathcal M_\beta
    -p^\mu F'p_\nu\sigma^\nu_{\alpha\dot\beta}
      \bar{\mathcal M}^{\dot\beta}\right]
  \\
  &+\langle X\lvert S^\mu_\alpha(0)
      \rvert Y\rangle_{\text{no pole}},
  \\
  \bra{X}\bar S^\mu_{\dot\alpha}(0)\ket{Y}
  ={}&\frac{\ii}{p^2}\left[
    F^*(\bar\sigma^\mu p_\nu\sigma^\nu)_{\dot\alpha}{}^{\dot\beta}
      \bar{\mathcal M}_{\dot\beta}
    -p^\mu F'{}^*p_\nu\bar\sigma^{\nu}_{\dot\alpha\beta}
      \mathcal M^\beta\right]
  \\
  &+\langle X\lvert\bar S^\mu_{\dot\alpha}(0)
      \rvert Y\rangle_{\text{no pole}}\,,
  \end{aligned}
  \tag{29.2.9}\label{eq:29.2.9}
\end{equation}
where \(\mathcal M_\alpha\) and
\(\bar{\mathcal M}_{\dot\alpha}\) are the undotted and dotted amplitudes
for emitting a goldstino with four-momentum \(p\), and the subscript
`no pole' denotes the terms in the matrix element that do not have the
one-goldstino pole in the four-momentum \(p\). Conservation of both current
components tells us that these matrix elements vanish if contracted with \(p_\mu\),
so in the limit \(p^\mu\to0\) the goldstino emission amplitude
is\footnote{The amplitude
\(\bra{X} S^\mu_\alpha(0)\ket{Y}_{\text{no pole}}\),
and likewise its dotted conjugate, may have
poles that go as \(1/p\mathbin{\cdot}k\) as \(p^\mu\to0\); they would not
arise from the goldstino propagator, which is explicitly excluded in this
matrix element, but rather from other particle propagators produced by the
insertion of the supersymmetry current in external lines of momentum \(k\)
of the process \(X\to Y\). In the limit \(p^\mu\to0\), the contribution
of these poles in Eq.~\eqref{eq:29.2.10} would dominate the amplitudes for
emission or absorption of soft goldstinos. But for such poles to arise,
the goldstino would have to be emitted in a transition between a pair of
degenerate particles of opposite statistics, which are not likely to
appear in theories in which supersymmetry is spontaneously broken. The
interactions of soft goldstinos differ in this respect from the
interactions of soft pions, photons, or gravitons.}
\begin{equation}
  \begin{aligned}
  \mathcal M_\alpha(X\to Y+g)
  &\xrightarrow[p^\mu\to0]{}
  -\frac{\ii}{F}\,
  p_\mu\langle X\lvert S^\mu_\alpha(0)
    \rvert Y\rangle_{\text{no pole}},
  \\
  \bar{\mathcal M}_{\dot\alpha}(X\to Y+g)
  &\xrightarrow[p^\mu\to0]{}
  -\frac{\ii}{F^*}\,
  p_\mu\langle X\lvert\bar S^\mu_{\dot\alpha}(0)
    \rvert Y\rangle_{\text{no pole}}\,.
  \end{aligned}
  \tag{29.2.10}\label{eq:29.2.10}
\end{equation}

There is another sum rule that provides an alternative proof of the
existence of goldstinos when supersymmetry is spontaneously broken, and
that also relates the parameter \(F\) and the vacuum energy density to the
\(D\)-terms and \(\mathcal{F}\)-terms that characterize the strength of
supersymmetry breaking. (This use of a sum rule is analogous to the second
proof in Section 19.2 of the existence of Goldstone bosons when ordinary
symmetries are spontaneously broken.) To derive this sum rule, we will now
give up the device of a finite volume, and instead avoid the question of
the convergence of the integral \eqref{eq:29.1.3} by working with a local
consequence of the supersymmetry of the action
\begin{equation}
  \left[
    \xi^\alpha S^0_\alpha(\mathbf x,t)
    +\bar\xi_{\dot\alpha}\bar S^{0\dot\alpha}(\mathbf x,t),
    \chi(\mathbf y,t)
  \right]
  =
  \ii\delta^3(\mathbf x-\mathbf y)\delta\chi(\mathbf x,t)+\cdots\,,
  \tag{29.2.11}\label{eq:29.2.11}
\end{equation}
where \(\chi(x)\) is an arbitrary fermionic or bosonic field,
\(\delta\chi(x)\) is the change in \(\chi(x)\) induced by the supersymmetry
transformation with infinitesimal parameters \(\xi_\alpha\) and
\(\bar\xi_{\dot\alpha}\) that leaves the
action invariant, and the dots denote terms involving derivatives of
\(\delta^3(\mathbf x-\mathbf y)\).

Let us consider the vacuum expectation values of the anticommutators of an
arbitrary undotted spinor field \(\psi_\alpha(x)\) with the two ordered
current components
\(\mathscr S^{\mu I}\equiv
(S^{\mu\beta},\bar S^\mu_{\dot\beta})\), where \(I\) records whether the
free current index is undotted or dotted. By summing over
a complete set \(\ket{X}\) of intermediate states (including
integration over particle momenta), this vacuum expectation value may be
put in the form
\begin{equation}
  \begin{aligned}
  &\langle\mathrm{VAC}\lvert
    \{\psi_\alpha(x),\mathscr S^{\mu I}(y)\}
    \rvert\mathrm{VAC}\rangle
  \\
  &\qquad
  =\int \dd^4p\,\exp\bigl(\ii p\mathbin{\cdot}(x-y)\bigr)
    \left[G^\mu_\alpha{}^I(p)
    +\widetilde G^\mu_\alpha{}^I(-p)\right]\,,
  \end{aligned}
  \tag{29.2.12}\label{eq:29.2.12}
\end{equation}
where
\begin{equation}
  G^\mu_\alpha{}^I(p)
  \equiv
  \sum_X\delta^4(p-p_X)
  \bra{\mathrm{VAC}}\psi_\alpha(0)\ket{X}
  \bra{X}\mathscr S^{\mu I}(0)\ket{\mathrm{VAC}}\,,
  \tag{29.2.13}\label{eq:29.2.13}
\end{equation}
\begin{equation}
  \widetilde G^\mu_\alpha{}^I(p)
  \equiv
  \sum_X\delta^4(p-p_X)
  \bra{\mathrm{VAC}}\mathscr S^{\mu I}(0)\ket{X}
  \bra{X}\psi_\alpha(0)\ket{\mathrm{VAC}}\,.
  \tag{29.2.14}\label{eq:29.2.14}
\end{equation}
Lorentz invariance requires that the undotted and dotted blocks of
\(G^\mu(p)\) and \(\widetilde G^\mu(p)\) take the forms
\begin{equation}
  \begin{aligned}
  G^\mu_{\alpha\dot\beta}(p)
  ={}&-\ii\theta(p^0)\left[
    \sigma^\mu_{\alpha\dot\beta}G^{(1)}(-p^2)
    +p_\nu\sigma^\nu_{\alpha\dot\beta}
      p^\mu G^{(2)}(-p^2)\right],
  \\
  G^\mu_\alpha{}^\beta(p)
  ={}&\theta(p^0)\left[
    p^\mu\delta_\alpha{}^\beta G^{(3)}(-p^2)
    -(p_\nu\sigma^\nu\bar\sigma^\mu)_\alpha{}^\beta
      G^{(4)}(-p^2)\right].
  \end{aligned}
  \tag{29.2.15}\label{eq:29.2.15}
\end{equation}
and
\begin{equation}
  \begin{aligned}
  \widetilde G^\mu_{\alpha\dot\beta}(p)
  ={}&-\ii\theta(p^0)\left[
    \sigma^\mu_{\alpha\dot\beta}\widetilde G^{(1)}(-p^2)
    +p_\nu\sigma^\nu_{\alpha\dot\beta}
      p^\mu\widetilde G^{(2)}(-p^2)\right],
  \\
  \widetilde G^\mu_\alpha{}^\beta(p)
  ={}&\theta(p^0)\left[
    p^\mu\delta_\alpha{}^\beta\widetilde G^{(3)}(-p^2)
    -(p_\nu\sigma^\nu\bar\sigma^\mu)_\alpha{}^\beta
      \widetilde G^{(4)}(-p^2)\right].
  \end{aligned}
  \tag{29.2.16}\label{eq:29.2.16}
\end{equation}
The right-hand sides of Eqs.~\eqref{eq:29.2.15} and
\eqref{eq:29.2.16} are unaffected if we replace \(-p^2\) with \(m^2\),
multiply by a factor \(\delta(p^2+m^2)\), and integrate over \(m^2\).
In this way, Eq.~\eqref{eq:29.2.12} becomes
\begin{equation}
  \begin{aligned}
  &(2\pi)^{-3}
  \langle\mathrm{VAC}\lvert
    \{\psi_\alpha(x),\bar S^\mu_{\dot\beta}(y)\}
  \rvert\mathrm{VAC}\rangle
  \\
  &=\ii\int_0^\infty \dd m^2
  \left[
    -\sigma^\mu_{\alpha\dot\beta}G^{(1)}(m^2)
    +\sigma^\nu_{\alpha\dot\beta}
      \partial_\nu\partial^\mu G^{(2)}(m^2)
  \right]\Delta_+(x-y,m)
  \\
  &\quad+\ii\int_0^\infty \dd m^2
  \left[
    -\sigma^\mu_{\alpha\dot\beta}\widetilde G^{(1)}(m^2)
    +\sigma^\nu_{\alpha\dot\beta}
      \partial_\nu\partial^\mu\widetilde G^{(2)}(m^2)
  \right]\Delta_+(y-x,m)\,,
  \\
  &(2\pi)^{-3}
  \langle\mathrm{VAC}\lvert
    \{\psi_\alpha(x),S^{\mu\beta}(y)\}
  \rvert\mathrm{VAC}\rangle
  \\
  &=\ii\int_0^\infty \dd m^2
  \left[
    -\delta_\alpha{}^\beta\partial^\mu G^{(3)}(m^2)
    +(\sigma^\nu\bar\sigma^\mu)_\alpha{}^\beta
      \partial_\nu G^{(4)}(m^2)
  \right]\Delta_+(x-y,m)
  \\
  &\quad+\ii\int_0^\infty \dd m^2
  \left[
    -\delta_\alpha{}^\beta\partial^\mu\widetilde G^{(3)}(m^2)
    +(\sigma^\nu\bar\sigma^\mu)_\alpha{}^\beta
      \partial_\nu\widetilde G^{(4)}(m^2)
  \right]\Delta_+(y-x,m)\,,
  \end{aligned}
  \tag{29.2.17}\label{eq:29.2.17}
\end{equation}
where \(\Delta_+(x,m)\) is the standard function
\begin{equation}
  \Delta_+(x,m)
  \equiv
  (2\pi)^{-3}\int \dd^4p\,
  \theta(p^0)\delta(p^2+m^2)\exp(\ii p\mathbin{\cdot}x)\,.
  \tag{29.2.18}\label{eq:29.2.18}
\end{equation}
Causality requires that the anticommutator on the left-hand side of
Eq.~\eqref{eq:29.2.17} should vanish for spacelike separations \(x-y\).
For such separations, \(\Delta_+(x-y,m)\) is an even function of \(x-y\),
so Eq.~\eqref{eq:29.2.17} vanishes for all spacelike separations if and
only if
\begin{equation}
  \begin{aligned}
  G^{(1)}(m^2)&=-\widetilde G^{(1)}(m^2)\,,
  &\qquad
  G^{(2)}(m^2)&=-\widetilde G^{(2)}(m^2)\,,
  \\
  G^{(3)}(m^2)&=+\widetilde G^{(3)}(m^2)\,,
  &
  G^{(4)}(m^2)&=+\widetilde G^{(4)}(m^2)\,,
  \end{aligned}
  \tag{29.2.19}\label{eq:29.2.19}
\end{equation}
so that for general \(x-y\), Eq.~\eqref{eq:29.2.17} now reads
\begin{equation}
  \begin{aligned}
  &(2\pi)^{-3}
  \langle\mathrm{VAC}\lvert
    \{\psi_\alpha(x),\bar S^\mu_{\dot\beta}(y)\}
  \rvert\mathrm{VAC}\rangle
  \\
  &\quad=\ii\int_0^\infty \dd m^2
  \left[
    -G^{(1)}(m^2)\sigma^\mu_{\alpha\dot\beta}
    +G^{(2)}(m^2)\sigma^\nu_{\alpha\dot\beta}
      \partial_\nu\partial^\mu
  \right]\Delta(x-y,m),
  \\
  &(2\pi)^{-3}
  \langle\mathrm{VAC}\lvert
    \{\psi_\alpha(x),S^{\mu\beta}(y)\}
  \rvert\mathrm{VAC}\rangle
  \\
  &\quad=\ii\int_0^\infty \dd m^2
  \left[
    -G^{(3)}(m^2)\delta_\alpha{}^\beta\partial^\mu
    +G^{(4)}(m^2)
      (\sigma^\nu\bar\sigma^\mu)_\alpha{}^\beta\partial_\nu
  \right]\Delta(x-y,m)\,,
  \end{aligned}
  \tag{29.2.20}\label{eq:29.2.20}
\end{equation}
where, as usual,
\begin{equation}
  \Delta(x-y,m)
  \equiv
  \Delta_+(x-y,m)-\Delta_+(y-x,m)\,.
  \tag{29.2.21}\label{eq:29.2.21}
\end{equation}
We next impose the supersymmetry current conservation condition
\eqref{eq:29.1.2}, which (because \(\Box\Delta=m^2\Delta\)) yields
\begin{equation}
  G^{(1)}(m^2)=m^2G^{(2)}(m^2)\,,
  \qquad
  m^2G^{(3)}(m^2)=-m^2G^{(4)}(m^2)\,.
  \tag{29.2.22}\label{eq:29.2.22}
\end{equation}
Finally, we must relate these spectral functions to the breaking of
supersymmetry. Recall that for \(x^0=y^0\),
\[
  \Delta(\mathbf x-\mathbf y,0,m)=0\,,
  \qquad
  \dot\Delta(\mathbf x-\mathbf y,0,m)
  =-\ii\delta^3(\mathbf x-\mathbf y)\,.
\]
Setting \(x^0=y^0=t\) and \(\mu=0\) in Eq.~\eqref{eq:29.2.20} thus yields
\begin{equation}
  \begin{aligned}
  &\langle\mathrm{VAC}\lvert
    \{\psi_\alpha(\mathbf x,t),S^{0\beta}(\mathbf y,t)\}
  \rvert\mathrm{VAC}\rangle
  \\
  &\quad=
  (2\pi)^3\delta^3(\mathbf x-\mathbf y)
  \delta_\alpha{}^\beta
  \int_0^\infty \dd m^2
  \left[G^{(3)}(m^2)+G^{(4)}(m^2)\right]\,.
  \end{aligned}
  \tag{29.2.23}\label{eq:29.2.23}
\end{equation}
Contracting on the right with an infinitesimal undotted supersymmetry
transformation parameter \(\xi_\beta\) and using
Eq.~\eqref{eq:29.2.11} then gives
\begin{equation}
  \ii\langle\delta\psi_\alpha\rangle_{\mathrm{VAC}}
  =(2\pi)^3\xi_\alpha
  \int_0^\infty \dd m^2
  \left[G^{(3)}(m^2)+G^{(4)}(m^2)\right]\,.
  \tag{29.2.24}\label{eq:29.2.24}
\end{equation}
But Eq.~\eqref{eq:29.2.22} shows that the integrand of the integral over
\(m^2\) in Eq.~\eqref{eq:29.2.24} vanishes, except perhaps at \(m^2=0\),
so we can conclude that
\begin{equation}
  G^{(3)}(m^2)+G^{(4)}(m^2)
  =\delta(m^2)\mathcal{G}\,,
  \tag{29.2.25}\label{eq:29.2.25}
\end{equation}
with a constant coefficient \(\mathcal{G}\) given by
\begin{equation}
  \ii\langle\delta\psi_\alpha\rangle_{\mathrm{VAC}}
  =(2\pi)^3\mathcal{G}\xi_\alpha\,.
  \tag{29.2.26}\label{eq:29.2.26}
\end{equation}

As we saw in Sections 26.4 and 27.4, a breakdown of supersymmetry is
signaled by the appearance of vacuum expectation values of the changes
\(\delta\psi\) under supersymmetry transformations of one or more spinor
fields \(\psi\). Eqs.~\eqref{eq:29.2.25} and \eqref{eq:29.2.26} show that
for any such spinor field, the spectral function
\(G^{(3)}(m^2)+G^{(4)}(m^2)\) has a delta function singularity at
\(m^2=0\), which could only arise from the appearance of a massless
one-particle state \(\ket{g}\) in the sums over states in
Eqs.~\eqref{eq:29.2.13} and/or \eqref{eq:29.2.14}. For the matrix elements
\(\bra{\mathrm{VAC}}\psi\ket{g}\) or
\(\bra{g}\psi\ket{\mathrm{VAC}}\) not to vanish, this
massless particle must have spin \(1/2\). This is the goldstino.

Now let's calculate the contribution of a one-goldstino state
\(\ket{\mathbf p,\lambda}\) with momentum \(\mathbf p\) and
helicity \(\lambda\) to the spectral functions \(G^{(i)}(m^2)\). Lorentz
invariance tells us that the matrix element of a general fermion field
(not just the renormalized goldstino field) in Eq.~\eqref{eq:29.2.13}
takes the form\footnote{Again, Lorentz invariance would allow independent
coefficients \(N_L\) and \(N_R\) for the undotted and dotted components,
respectively. Using the CPT transformation of the
one-particle state and the reality property of the coefficient functions
discussed in the first footnote of this section, and the CPT
transformation of a general fermion field,
\(\mathrm{CPT}\,\psi_\alpha(x)(\mathrm{CPT})^{-1}
=\bar\psi_{\dot\alpha}(-x)\), we find that \(N_L=N_R^*\).}
\begin{equation}
  \begin{aligned}
  \langle\mathrm{VAC}\lvert\psi_\alpha(0)
    \rvert\mathbf p,\lambda\rangle
  &=\frac{N}{(2\pi)^{3/2}}u_\alpha(\mathbf p,\lambda),
  \\
  \langle\mathrm{VAC}\lvert\bar\psi_{\dot\alpha}(0)
    \rvert\mathbf p,\lambda\rangle
  &=\frac{N^*}{(2\pi)^{3/2}}
    \bar u_{\dot\alpha}(\mathbf p,\lambda),
  \end{aligned}
  \tag{29.2.27}\label{eq:29.2.27}
\end{equation}
where \(N\) is a constant, characterizing the particular fermion field.
Also, the delta function in Eq.~\eqref{eq:29.2.13} is
\[
  \delta^4(p-p_g)
  =2p^0\delta(p^2)\theta(p^0)
    \delta^3(\mathbf p-\mathbf p_g)\,.
\]
Together with Eq.~\eqref{eq:29.2.7}, this yields
\[
  \begin{aligned}
  \left[G^\mu_{\alpha\dot\beta}(p)\right]_1
  &=
  \frac{NF'}{(2\pi)^3}\theta(p^0)\delta(p^2)
    p^\mu p_\nu\sigma^\nu_{\alpha\dot\beta},
  \\
  \left[G^\mu_\alpha{}^\beta(p)\right]_1
  &=
  -\frac{NF}{(2\pi)^3}\theta(p^0)\delta(p^2)
    (p_\nu\sigma^\nu\bar\sigma^\mu)_\alpha{}^\beta .
  \end{aligned}
\]
The subscript 1 indicates the one-goldstino contribution. Comparison with
Eq.~\eqref{eq:29.2.15} shows that
\begin{equation}
  \begin{aligned}
  \left[G^{(2)}(m^2)\right]_1
  &=\ii(2\pi)^{-3}NF'\delta(m^2)\,,
  &\qquad
  \left[G^{(4)}(m^2)\right]_1
  &=(2\pi)^{-3}NF\delta(m^2)\,,
  \\
  \left[G^{(1)}(m^2)\right]_1
  &=\left[G^{(3)}(m^2)\right]_1=0\,.
  \end{aligned}
  \tag{29.2.28}\label{eq:29.2.28}
\end{equation}
Thus Eqs.~\eqref{eq:29.2.25} and \eqref{eq:29.2.26} yield
\begin{equation}
  \ii\langle\delta\psi_\alpha\rangle_{\mathrm{VAC}}
  =NF\xi_\alpha\,.
  \tag{29.2.29}\label{eq:29.2.29}
\end{equation}
Evidently if the change in any fermion field under a supersymmetry
transformation has a non-zero vacuum expectation value then the
\(N\)-factor for that field cannot vanish, so there must be a
one-goldstino state that contributes to these spectral functions.

To be a little more specific, recall that the fermionic components
\(\psi_{n\alpha}\) of a left-chiral scalar superfield \(\Phi_n\) obey the
supersymmetry transformation rule \eqref{eq:26.3.15}:
\begin{equation}
  \delta\psi_{n\alpha}
  =-\ii\sqrt{2}\,
    \sigma^\mu_{\alpha\dot\alpha}\bar\xi^{\dot\alpha}
      \partial_\mu\phi_n
  +\sqrt{2}\mathcal{F}_n\xi_\alpha\,,
  \tag{29.2.30}\label{eq:29.2.30}
\end{equation}
while Eq.~\eqref{eq:27.3.5} gives the supersymmetry transformation of the
gaugino fields as
\begin{equation}
  \delta\lambda_{A\alpha}
  =f_{A\mu\nu}(\sigma^{\mu\nu}\xi)_\alpha
    +\ii D_A\xi_\alpha\,.
  \tag{29.2.31}\label{eq:29.2.31}
\end{equation}
Hence the \(N\)-factors in the matrix elements of the \(\psi_n\) and
\(\lambda_A\) between the vacuum and a one-goldstino state are given by
\begin{equation}
  N_n
  =\ii F^{-1}\sqrt{2}
    \langle\mathcal{F}_n\rangle_{\mathrm{VAC}}\,,
  \qquad
  N_A
  =-F^{-1}\langle D_A\rangle_{\mathrm{VAC}}\,.
  \tag{29.2.32}\label{eq:29.2.32}
\end{equation}

Let us see how the results \eqref{eq:29.2.32} arise in the tree
approximation. For a renormalizable theory of gauge and chiral
superfields, Eq.~\eqref{eq:27.4.30} gives the fermion mass matrix in an
undotted Weyl basis as
\begin{equation}
  \begin{aligned}
  M_{nm}
  &=
  \left(
    \frac{\partial^2f(\phi)}
         {\partial\phi_n\partial\phi_m}
  \right)_{\phi=\phi_0}\,,
  \\
  M_{nA}=M_{An}
  &=\ii\sqrt{2}(t_A\phi_0)_n^*\,,
  \qquad
  M_{AB}=0\,.
  \end{aligned}
  \tag{29.2.33}\label{eq:29.2.33}
\end{equation}
The vacuum fields \(\phi_{n0}\) are at a minimum of the potential given
by Eq.~\eqref{eq:27.4.9}, so
\begin{equation}
  0
  =
  \left(\frac{\partial V(\phi)}{\partial\phi_n}\right)_0
  =
  -\sum_m M_{nm}\mathcal{F}_{m0}
  +\sum_A(\phi_0^\dagger t_A)_nD_{A0}\,,
  \tag{29.2.34}\label{eq:29.2.34}
\end{equation}
where the \(\mathcal{F}\)s and \(D\)s are given by
Eqs.~\eqref{eq:27.4.6} and \eqref{eq:27.4.7}
\[
  \mathcal{F}_n
  =-\left(\partial f(\phi)/\partial\phi_n\right)^*\,,
  \qquad
  D_A
  =\xi_A+\sum_{nm}\phi_n^*(t_A)_{nm}\phi_m\,,
\]
and the subscript 0 indicates that we are to set
\(\phi_n=\phi_{n0}\). Furthermore, the gauge invariance of the
superpotential requires that
\begin{equation}
  \sum_n\mathcal{F}_n(t_B\phi)_n^*=0\,,
  \tag{29.2.35}\label{eq:29.2.35}
\end{equation}
for all values of \(\phi\). Therefore the undotted fermion mass matrix
\(M\) has an eigenvector \(v\) with \(Mv=0\), where
\begin{equation}
  v_n=\sqrt{2}\mathcal{F}_{n0}\,,
  \qquad
  v_A=\ii D_{A0}\,.
  \tag{29.2.36}\label{eq:29.2.36}
\end{equation}
Thus in the tree approximation, if we expand the undotted fermion
fields in renormalized fields for particles of definite mass, then the
coefficients of the goldstino field in \(\psi_{n\alpha}\) and in
\(\lambda_{A\alpha}\) are proportional to \(\sqrt{2}\mathcal{F}_{n0}\) and
\(iD_{A0}\), respectively, in agreement with Eq.~\eqref{eq:29.2.32}.

If we normalize the spinor fields \(\psi_{n\alpha}\) and
\(\lambda_{A\alpha}\) so that
the matrix connecting these fields to the renormalized fields of particles
of definite mass is unitary, then
\begin{equation}
  \sum_n\lvert N_n\rvert^2+\sum_A\lvert N_A\rvert^2=1\,,
  \tag{29.2.37}\label{eq:29.2.37}
\end{equation}
so Eq.~\eqref{eq:29.2.32} gives the non-perturbative result
\begin{equation}
  \lvert F\rvert^2
  =
  2\sum_n
  \left\lvert\langle\mathcal{F}_n\rangle_{\mathrm{VAC}}\right\rvert^2
  +\sum_A
  \left\lvert\langle D_A\rangle_{\mathrm{VAC}}\right\rvert^2\,.
  \tag{29.2.38}\label{eq:29.2.38}
\end{equation}
The result \eqref{eq:29.2.38} for \(\lvert F\rvert^2\) allows us to
express the vacuum expectation density \eqref{eq:29.2.8} in terms of the
vacuum expectation values of auxiliary fields
\begin{equation}
  \rho_{\mathrm{VAC}}
  =
  \sum_n
  \left\lvert\langle\mathcal{F}_n\rangle_{\mathrm{VAC}}\right\rvert^2
  +\frac12\sum_A
  \left\lvert\langle D_A\rangle_{\mathrm{VAC}}\right\rvert^2\,.
  \tag{29.2.39}\label{eq:29.2.39}
\end{equation}
This is the non-perturbative generalization of the zeroth-order result
\eqref{eq:27.4.9}. It confirms a result used in Section 27.6, that the
conditions \(\langle\mathcal{F}_n\rangle=\langle D_A\rangle=0\) are
sufficient as well as necessary for supersymmetry to be unbroken.

\begin{center}
\(\ast\qquad\ast\qquad\ast\)
\end{center}

It is instructive to see how Eq.~\eqref{eq:29.2.8} can be derived without
the device of a finite volume. For this purpose, consider the vacuum
expectation value of the anticommutator of two supersymmetry currents. By
using Lorentz invariance and the vanishing of anticommutators at spacelike
separations in the same way as in the previous section, we find
\begin{equation}
  \begin{aligned}
  &\langle\mathrm{VAC}\lvert
    \{S^\mu_\alpha(x),\bar S^\nu_{\dot\beta}(y)\}
  \rvert\mathrm{VAC}\rangle
  \\
  &\quad
  =-\int \dd m^2
  \bigl[
    H^{(1)}(m^2)\sigma^\mu_{\alpha\dot\beta}\partial^\nu
    +H^{(2)}(m^2)\sigma^\nu_{\alpha\dot\beta}\partial^\mu
    +H^{(3)}(m^2)\sigma^\rho_{\alpha\dot\beta}
      \partial_\rho\partial^\mu\partial^\nu
  \\
  &\hspace{38mm}
    +H^{(4)}(m^2)\sigma^\rho_{\alpha\dot\beta}
      \partial_\rho\eta^{\mu\nu}
    +H^{(5)}(m^2)\epsilon^{\mu\nu\lambda\rho}
      \partial_\lambda\sigma_{\rho\,\alpha\dot\beta}
  \bigr]\Delta(x-y,m)
  +\cdots\,.
  \end{aligned}
  \tag{29.2.40}\label{eq:29.2.40}
\end{equation}
Here
\begin{equation}
  \begin{aligned}
  &\int \dd X\,\delta^4(p-p_X)
  \bra{\mathrm{VAC}} S^\mu_\alpha(0)\ket{X}
  \bra{X}\bar S^\nu_{\dot\beta}(0)\ket{\mathrm{VAC}}
  \\
  &\quad
  =-\ii\bigl[
  H^{(1)}(-p^2)\sigma^\mu_{\alpha\dot\beta}p^\nu
  +H^{(2)}(-p^2)\sigma^\nu_{\alpha\dot\beta}p^\mu
  -H^{(3)}(-p^2)p_\rho\sigma^\rho_{\alpha\dot\beta}p^\mu p^\nu
  +H^{(4)}(-p^2)p_\rho\sigma^\rho_{\alpha\dot\beta}\eta^{\mu\nu}
  \\
  &\hspace{29mm}
  +H^{(5)}(-p^2)\epsilon^{\mu\nu\lambda\rho}
    p_\lambda\sigma_{\rho\,\alpha\dot\beta}\bigr]+\cdots\,,
  \end{aligned}
  \tag{29.2.41}\label{eq:29.2.41}
\end{equation}
and the dots in Eqs.~\eqref{eq:29.2.40} and \eqref{eq:29.2.41} denote the
independent same-chirality scalar and antisymmetric-tensor spinor
structures, which will not concern us here. The conjugacy of the undotted
and dotted currents together with Eq.~\eqref{eq:26.A.20} tells us that
the spectral functions \(H^{(i)}\) are all real, and the
conservation of the supersymmetry currents dictate that
\begin{equation}
  H^{(1)}(m^2)
  =H^{(2)}(m^2)
  =-m^2H^{(3)}(m^2)-H^{(4)}(m^2)\,,
  \tag{29.2.42}\label{eq:29.2.42}
\end{equation}
and
\begin{equation}
  m^2H^{(1)}(m^2)=0\,.
  \tag{29.2.43}\label{eq:29.2.43}
\end{equation}
Setting \(\mu=\nu=0\) and \(x^0=y^0\) in Eq.~\eqref{eq:29.2.40},
integrating over \(\mathbf x\), and using Eq.~\eqref{eq:29.2.42} then
gives
\begin{equation}
  \begin{aligned}
  \langle\mathrm{VAC}\lvert
  \{Q_\alpha,\bar S^0_{\dot\beta}(0)\}
  \rvert\mathrm{VAC}\rangle
  ={}&(2\pi)^3\sigma^0_{\alpha\dot\beta}\int \dd m^2
  \bigl[
    H^{(1)}(m^2)+H^{(2)}(m^2)
  \\
  &\qquad
    +m^2H^{(3)}(m^2)+H^{(4)}(m^2)
  \bigr]
  +\cdots
  \\
  ={}&(2\pi)^3\sigma^0_{\alpha\dot\beta}
    \int \dd m^2H^{(1)}(m^2)+\cdots\,.
  \end{aligned}
  \tag{29.2.44}\label{eq:29.2.44}
\end{equation}
To evaluate this anticommutator, we first write it as
\begin{equation}
  \{Q_\alpha,\bar S^\nu_{\dot\beta}(x)\}
  =-2\sigma_{\mu\,\alpha\dot\beta}T^{\mu\nu}(x)+\cdots\,,
  \tag{29.2.45}\label{eq:29.2.45}
\end{equation}
where \(T^{\mu\nu}(x)\) is some tensor operator and the dots again denote
the complementary same-chirality spinor structures. Hermitian conjugacy
of \(Q_\alpha,\bar Q_{\dot\alpha}\) and of the two currents tells us that
\(T^{\mu\nu}(x)\) is Hermitian; conservation of the supersymmetry current tells us that
it is conserved, in the sense that
\begin{equation}
  \partial_\nu T^{\mu\nu}=0\,,
  \tag{29.2.46}\label{eq:29.2.46}
\end{equation}
and the anticommutation relation \eqref{eq:25.2.36} tells us that
\begin{equation}
  \int \dd^3x\,T^{\mu0}(x)=P^\mu\,.
  \tag{29.2.47}\label{eq:29.2.47}
\end{equation}
These properties allow us to identify \(T^{\mu\nu}\) as the
energy-momentum tensor. It is not in general equal to the symmetric
energy-momentum tensor \(\Theta^{\mu\nu}\) discussed in Section 7.4, but
Eq.~\eqref{eq:29.2.47} shows that the energy density \(T^{00}\) can
differ from \(\Theta^{00}\) only by spatial derivative terms that cannot
contribute in states of zero three-momentum, so the vacuum energy density
is given by
\begin{equation}
  \rho_{\mathrm{VAC}}
  =\bra{\mathrm{VAC}} T^{00}\ket{\mathrm{VAC}}\,.
  \tag{29.2.48}\label{eq:29.2.48}
\end{equation}
Thus Eqs.~\eqref{eq:29.2.44} and \eqref{eq:29.2.45} yield
\begin{equation}
  2\rho_{\mathrm{VAC}}
  =(2\pi)^3\int \dd m^2H^{(1)}(m^2)\,.
  \tag{29.2.49}\label{eq:29.2.49}
\end{equation}
But Eq.~\eqref{eq:29.2.43} tells us that \(H^{(1)}(m^2)\) vanishes
except perhaps at \(m^2=0\), so
\begin{equation}
  H^{(1)}(m^2)
  =2(2\pi)^{-3}\delta(m^2)\rho_{\mathrm{VAC}}\,.
  \tag{29.2.50}\label{eq:29.2.50}
\end{equation}
Thus we see again that a non-vanishing vacuum energy density entails the
existence of a massless fermion, the goldstino. Using Eq.~\eqref{eq:29.2.7},
a straightforward calculation of the one-goldstino contribution to the
spectral functions yields
\begin{equation}
  H^{(1)}(m^2)
  =(2\pi)^{-3}\delta(m^2)\lvert F\rvert^2\,.
  \tag{29.2.51}\label{eq:29.2.51}
\end{equation}
Comparing this with Eq.~\eqref{eq:29.2.50} gives our previous result
\eqref{eq:29.2.8} for the vacuum energy density.
